% P05.2 one-dimensional graph classification, 25 September 2026.

Continue with the finite one-dimensional balanced rational polyhedral
structure $\Pi$ and the local relation cones $K_v$ above.  We may insert one
bivalent vertex in every component without a vertex.  This does not change
any original relation cone or the orthogonality problem.

\begin{definition}
\label{def:frontier-polyhedral-mass-accessible}
Let $\gamma\in\mathrm{PA}_{\mathbf R}(\mathbf X)\cap
\mathrm{PC}(\mathbf X)$.  A vertex $v\in\Pi^{(0)}$ is
$\gamma$-mass-accessible if there is a function
\[
 \varphi\in\mathrm{PA}_{\mathbf R}(\mathbf X)\cap
 \mathrm{PC}(\mathbf X),
 \qquad
 \varphi-\gamma\in\mathrm{PA}_{\mathbf R,b}(\mathbf X),
\]
such that
\[
 \MA_{\mathrm{poly},\mathbf R}(\varphi)\{v\}
 =\sum_{e\in E(v)}w_es_{v,e}(\varphi)>0.
\]
\end{definition}

\begin{lemma}
\label{lem:frontier-polyhedral-ambient-vertex-accessibility}
Let $\gamma$ be positive at infinity.  Then every vertex of $\Pi$ is
$\gamma$-mass-accessible.
\end{lemma}

\begin{lemma}
\label{lem:frontier-polyhedral-accessible-mass-obstruction}
Suppose that $\dim K_v>1$ and that $v$ is
$\gamma$-mass-accessible.  Then $(\mathbf X,\gamma)$ does not have
real-polyhedral orthogonality.  More precisely, there is a bounded
real-piecewise-affine obstacle $u$ such that
\[
 \int_{\mathbf X}
 \bigl(P_\gamma(u)-(\gamma+u)\bigr)
 \MA_{\mathrm{poly},\mathbf R}(P_\gamma(u))<0.
\]
\end{lemma}

\begin{theorem}
\label{thm:frontier-polyhedral-graph-orthogonality-classification}
Let $\mathbf X$ be a one-dimensional balanced rational polyhedral space
admitting a finite structure $\Pi$, and let
$\gamma\in\mathrm{PA}_{\mathbf R}(\mathbf X)\cap
\mathrm{PC}(\mathbf X)$ be positive at infinity.  Then
$(\mathbf X,\gamma)$ has real-polyhedral orthogonality if and only if
\[
                         \dim K_v=1
                         \qquad\text{for every }v\in\Pi^{(0)}.
\]
Thus the repaired one-dimensional classification is local at the vertices:
bounded edges, graph cycles, and lattice saturation introduce no additional
condition.
\end{theorem}

Positivity at infinity enters the converse only through
\Cref{lem:frontier-polyhedral-ambient-vertex-accessibility}: it permits an
ambient polyhedral norm with strictly smaller recession slope on every
unbounded edge.  The statement neither repairs the literal integral-slope
class nor addresses higher-dimensional balanced spaces.
